\documentclass{article}

\usepackage[utf8]{inputenc}     % pdfLaTeX
\usepackage[T1]{fontenc}
\usepackage[magyar]{babel}
\usepackage{amsmath}

% Set page size and margins
% Replace `letterpaper' with `a4paper' for UK/EU standard size
\usepackage[a4paper,top=2cm,bottom=2cm,left=3cm,right=3cm,marginparwidth=1.75cm]{geometry}
%for arrows
\usepackage{textcomp}

\title{Intézménynevek Helyesírása}
\author{Lengyel Zoltán Péter \thanks{Csukás Ibolya tanítása alapján}}
\date{Legutóbb frissítve: 2026. Március 15.}

\begin{document}

\maketitle

\tableofcontents

\newpage

\section{Állandó címek}
\begin{itemize}
    \item Folyóiratok, rendszeres kiadványok 
    \item Minden lényeges szavuk nagybetű
    \begin{itemize}
        \item Nemzeti Sport
        \item Élet és Tudomány
        \item Élet és Irodalom
        \item Sport plusz Foci
        \item 168 Óra
    \end{itemize}
    \item toldalékolás: közvetlenül
    \begin{itemize}
        \item Élet és Tudományt
    \end{itemize}
    \item -i / -beli / -féle / -fajta toldalék: kötőjellel
    \begin{itemize}
        \item pl. Nemzeti Sport-beli cikkek
    \end{itemize}
    \item összetételi tag kapcsolása: kötőjellel
    \begin{itemize}
        \item pl. Nemzeti Sport-cikkek
    \end{itemize}
\end{itemize}

\section{Egyedi címek}
\begin{itemize}
    \item Csak az első szót és a tulajdonnev(ek)et írjuk nagybetűvel
    \begin{itemize}
        \item A Pál utcai fiúk
        \item Egri csillagok
    \end{itemize}
    \item Toldalékolás: közvetlenül
    \begin{itemize}
        \item a Négy évszakot
        \item Arany János Toldija
    \end{itemize}
    \item  -i / -beli / -féle / -fajta toldalék: kötőjellel
    \begin{itemize}
        \item Toldi-féle
        \item Négy évszak-beli
    \end{itemize}
    \item Ha a cím ragra végződik: kötőjellel kapcsoljuk az egyéb toldalékot
    \begin{itemize}
        \item Szülőföldemen-t
    \end{itemize}
    \item Ha a cím írásjelre végződik: kötőjellel kapcsoljuk a toldalékot
    \begin{itemize}
        \item pl. a Szerelem?-ből olvas
    \end{itemize}
\end{itemize}

\end{document}